% !Mode:: "TeX:UTF-8"
% The main file for the thesis.
% main.tex
%

% 编译方式,可选的有xelatex, pdflatex, dvipdfmx, dvipspdf, yap
% 推荐使用xelatex，默认UTF8编码格式，对中文支持很好。
\def\usewhat{xelatex}

% 若用xelatex编译，默认用UTF8编码，非UTF8文件时，需在每个文件中指定字符编码;
% 本段无需改动，需放在输入中文之前；
\def\atempxetex{xelatex}\ifx\atempxetex\usewhat
% 若文件使用非utf8编码,请将注释去掉
% \XeTeXinputencoding "gbk"
\fi

% input "reference\reference.bib" %for winedt users
% 版本号: X.Y.Z, X: major, Y: minor, Z: patch
\def\version{1.0.0}     % 该变量仅用于模板文件的版本号控制

\def\xuewei{Doctor}   % 定义学位 博士
%\def \xuewei {Master}  % 硕士

\def\oneortwoside{twoside} %定义单双面打印，只对硕士学位论文有效；
%\def\oneortwoside{oneside} % 硕士单面打印

\def\xueke{Engineering} % 定义学科 工学
%\def\xueke{Science}      % 理学
%\def\xueke{Management}   % 管理学
%\def\xueke{Arts}         % 艺术学

% dicta 2009 definition.
\DeclareMathAlphabet{\mathpzc}{OT1}{pzc}{m}{it}
\def\calV{\mathcal{V}}
\def\calE{\mathcal{E}}
\def\calC{\mathcal{C}}
\def\calU{\mathcal{U}}
\def\bfX{\mathbf{X}}
\def\bfx{\mathbf{x}}
\def\bfY{\mathbf{Y}}
\def\bfy{\mathbf{y}}
\def\bfL{\mathbf{L}}
\def\bfE{\mathbf{E}}
\def\bfD{\mathbf{D}}
\def\bfW{\mathbf{W}}
\def\bfT{\mathpzc{T}}
\def\bfone{\mathbf{1}}
\def\overX{\overline{X}}
\def\vector{\mathrm{vec}}
\def\lsep{\!\!&}
\def\rsep{&\!\!}
% end of dicta 2009

% !Mode:: "TeX:UTF-8"
% 若xelatex编译非UTF8文件，需在每个文件中指定字符编码;
% main.tex中手动制定了\atemp和\usewhat参数；
\ifx\atempxetex\usewhat
%\XeTeXinputencoding "gbk"
\fi

% 硕博类型 的一些定义

% 导言区使用中文
\makeatletter
\@tempcnta=128
\loop \catcode\@tempcnta=13 \ifnum\@tempcnta<255 \advance \@tempcnta \@ne
\repeat
\makeatother

\newif\ifxueweidoctor %判断论文类型
\newif\ifxueweimaster
\def\temp{Doctor}
\ifx\temp\xuewei
  \xueweidoctortrue  \xueweimasterfalse
\fi
\def\temp{Master}
\ifx\temp\xuewei
  \xueweidoctorfalse  \xueweimastertrue
\fi

\ifxueweidoctor
  \newcommand{\cxuewei}{博士}
  \newcommand{\exuewei}{Doctor}
  \newcommand{\exueweier}{Doctoral}
  \newcommand{\xueweishort}{博}
\fi

\ifxueweimaster
  \newcommand{\cxuewei}{硕士}
  \newcommand{\exuewei}{Master}
  \newcommand{\exueweier}{Master}
  \newcommand{\xueweishort}{硕}
\fi


\ifxueweidoctor
  \def\oneortwoside{twoside}
\fi

\ifx\oneortwoside\undefined
  \def\oneortwoside{twoside}
\fi

\newif\ifoneortwoside
\def\temp{twoside}
\ifx\temp\oneortwoside
  \oneortwosidetrue
\else
  \oneortwosidefalse
\fi
    % 硕博类型

%下面的book选项中可以使用 draft 选项，使插入的图形只显示外框，以加快预览速度。
\documentclass[12pt,a4paper,openany,\oneortwoside]{book}
% !Mode:: "TeX:UTF-8"
% 若xelatex编译非UTF8文件，需在每个文件中指定字符编码;
% main.tex中手动制定了\atemp和\usewhat参数；
\ifx\atempxetex\usewhat
%\XeTeXinputencoding "gbk"
\fi

\usepackage{dsfont}        % add by Tina 数学字体
%\usepackage{algorithm}     % add by Tina 写算法
%\usepackage{algorithmic}   % add by Tina 写算法
\usepackage{amssymb}
\usepackage{amsmath}        % AMSLaTeX宏包 用来排出更加漂亮的公式
%\usepackage{amsthm}

% 图形支持宏包 为了使用pdftex 需要作相应判断
\usepackage{etex}%增加计数器总数（原来是256，宏包多，可能不够用），编译需基于eTeX
\usepackage{ifpdf}
%定义一个新判断命令 %有一个宏包ifpdf 可以完成这件事，应该比这个严谨
%\newif\ifpdf
%\ifx\pdfoutput\undefined
%   \pdffalse
%\else
%   \pdfoutput=1
%   \pdftrue
%\fi
\usepackage{pifont}%调用文本个带圈的宏包
%%%%%%%%彩色引用和书签%%%%%%
\ifx\atempxetex\usewhat
    \usepackage[dvipdfm]{graphicx}
\else
    \ifpdf
        \usepackage[pdftex]{graphicx}
    \else
        \usepackage[dvips]{graphicx}
    \fi
\fi
\usepackage[top=30mm,bottom=24mm,left=25mm,right=25mm]{geometry}
\usepackage{layouts}                    % 打印当前页面格式的宏包
\usepackage[sf]{titlesec}               % 控制标题的宏包
\usepackage{titletoc}                   % 控制目录的宏包
\usepackage[perpage,symbol]{footmisc}   % 脚注控制
\usepackage{fancyhdr}                   % fancyhdr宏包 页眉和页脚的相关定义
\usepackage{fancyref}
\usepackage{array}          %增强表格的功能 %可能与xeCJK宏包冲突，需放在xeCJK之前。

\ifx\atempxetex\usewhat
\usepackage[slantfont,boldfont,CJKaddspaces]{xeCJK} %
%\CJKlanguage{zh-cn}
\else
\usepackage{CJK,CJKpunct}   % 中文支持宏包
\usepackage{times}          % 使用Times字体的宏包
\fi

\usepackage{type1cm}        % tex1cm宏包，控制字体的大小
\usepackage{indentfirst}    % 首行缩进宏包
\usepackage{color}          % 支持彩色

\usepackage{relsize}            % 调整公式字体大小 \mathsmaller \mathlarger
\usepackage{textcomp}         % 千分号等特殊符号
\usepackage{mathrsfs}       % 不同于\mathcal or \mathfrak 之类的英文花体字体
\usepackage{bm}              % 处理数学公式中的黑斜体的宏包
\usepackage[amsmath,thmmarks,hyperref]{ntheorem}% 定理类环境宏包，其中 amsmath 选项用来兼容 AMS LaTeX 的宏包

\usepackage{epsfig}         % eps图像
\usepackage[below]{placeins}%允许上一个section的浮动图形出现在下一个section的开始部分,还提供\FloatBarrier命令,使所有未处理的浮动图形立即被处理
%\usepackage{psfrag}        %替换eps图形中的文字
%\usepackage{floatflt}       % 图文混排用宏包
\usepackage{rotating}       % 图形和表格的控制
%\usepackage{endfloat}      %可将浮动对象放置到文件的最后
\usepackage{setspace}       % 定制表格和图形的多行标题行距
\usepackage{flafter}       % 使得所有浮动体不能被放置在其浮动环境之前，以免浮动体在引述它的文本之前出现.
\usepackage{multirow}       %使用Multirow宏包，使得表格可以合并多个row格
\usepackage{booktabs}       % 表格，横的粗线；\specialrule{1pt}{0pt}{0pt}
\usepackage{longtable}      %支持跨页的表格。
%\usepackage[centerlast]{caption2}       %浮动图形和表格标题样式，这里已被ccaption完全替代。
\usepackage[hang,center]{subfigure}%支持子图 %centerlast 设置最后一行是否居中
\usepackage[subfigure]{ccaption} %,caption2,

%\usepackage{cite}          % 支持引用的宏包
\usepackage[sort&compress,numbers]{natbib}% 支持引用缩写的宏包
%\usepackage{hypernat}

\usepackage{enumitem}       %使用enumitem宏包,改变列表项的格式
\usepackage{calc}           %长度可以用+ - * / 进行计算
\usepackage[boxed,linesnumbered,algochapter]{algorithm2e}    % 算法的宏包

% 生成有书签的pdf及其开关, 该宏包应放在所有宏包的最后, 宏包之间有冲突
\def\atemp{dvipspdf}\ifx\atemp\usewhat
\usepackage[dvips,
            CJKbookmarks=true,
            bookmarksnumbered=true,
            bookmarksopen=true,
            colorlinks=false,    % 最后打印的时候可以改成false，这样字体都是黑色
            pdfborder={0 0 1},   % 去掉链接的边框
            citecolor=blue,
            linkcolor=red,
            anchorcolor=green,
            urlcolor=blue,
            unicode,
            breaklinks=true
            ]{hyperref}
\usepackage{breakurl} %消除dvips的时候网页链接断行失效的问题。
\fi

\def\atemp{dvipdfmx}\ifx\atemp\usewhat
\usepackage[dvipdfm, %dvi-->pdf 生成书签
            CJKbookmarks=true,
            bookmarksnumbered=true,
            bookmarksopen=true,
            colorlinks=false,
            pdfborder={0 0 1},
            citecolor=blue,
            linkcolor=red,
            anchorcolor=green,
            urlcolor=blue,
            breaklinks=true
            ]{hyperref}
\AtBeginDvi{\special{pdf:tounicode GBK-EUC-UCS2}} % GBK -> Unicode,此时可以不用gbk2uni
\fi

\def\atemp{pdflatex}\ifx\atemp\usewhat
\usepackage{cmap}                       %pdflatex编译时，可以生成可复制、粘贴的中文PDF文档
\usepackage[pdftex,
            CJKbookmarks=true,
            bookmarksnumbered=true,
            bookmarksopen=true,
            colorlinks=false,
            pdfborder={0 0 1},
            citecolor=blue,
            linkcolor=red,
            anchorcolor=green,
            urlcolor=blue,
            unicode,
            breaklinks=true
            ]{hyperref}
\fi

\def\atemp{yap}\ifx\atemp\usewhat
\usepackage[dvipdf,  %考虑到yap可以正反向搜索，dvipdf使yap中的链接有效。
            CJKbookmarks=true,
            bookmarksnumbered=true,
            bookmarksopen=true,
            colorlinks=false,
            pdfborder={0 0 1},
            citecolor=blue,
            linkcolor=red,
            anchorcolor=green,
            urlcolor=blue,
            unicode,
            breaklinks=true
            ]{hyperref}
\fi

\ifx\atempxetex\usewhat %\def\atempxetex{xelatex} main.tex中已定义;
\usepackage[xetex,
            CJKbookmarks=true,
            bookmarksnumbered=true,
            bookmarksopen=true,
            colorlinks=false,
            pdfborder={0 0 1},
            citecolor=blue,
            linkcolor=red,
            anchorcolor=green,
            urlcolor=blue,
            breaklinks=true,
            naturalnames  %与algorithm2e宏包协调
            ]{hyperref}
\fi

\usepackage{arydshln}       %分块矩阵画虚线，挺好用

%\usepackage[hang]{caption}  % 标题的其它属性的自由控制 --- added by Tina
 % 引用的宏包

% 论文包含的内容
\includeonly{
            appendix/Authorization,
            body/chapter1,
            body/chapter2,
            body/chapter3,
            body/chapter4,
            body/chapter5,
            body/conclusion,
            body/acknowledgements,
            %appendix/appA,
            appendix/publications,  % 编译此文件时 一定要配合一章内容 否则会报错 --- by Tina
%            appendix/Resume
            }
\graphicspath{{figures/}} % 定义所有的eps文件在 figures 子目录下

\begin{document}
\ifx\atempxetex\usewhat\else
\begin{CJK*}{UTF8}{song}
\fi

% !Mode:: "TeX:UTF-8"
% 若xelatex编译非UTF8文件，需在每个文件中指定字符编码;
% main.tex中手动制定了\atemp和\usewhat参数；
\ifx\atempxetex\usewhat
%\XeTeXinputencoding "gbk"
\fi

%%%%%%%%%%%%%%%%%%%%%%%%%%%%%%%%%%%%%%%%%%%%%%%%%%%%%%%%%%%
% 重定义字体命令
% 注意win2000,没有 simsun,　最好到网上找一个
% 一些字体是office2000 带的
%%%%%%%%%%%%%%%%%%%%%%%%%%%%%%%%%%%%%%%%%%%%%%%%%%%%%%%%%%%
\ifx\atempxetex\usewhat %xelatex调用系统字体；
\setmainfont{Times New Roman}   % 正文中的英文 采用 Times New Roman 字体；
\setsansfont{SimHei}   % 正文中的英文 采用 Times New Roman 字体；
\setCJKmainfont{SimSun}
\setCJKsansfont{SimHei}
\setCJKmonofont{LiSu}
\setCJKfamilyfont{song}{SimSun}
\setCJKfamilyfont{hei}{SimHei}
\setCJKfamilyfont{fs}{FangSong}
\setCJKfamilyfont{kai}{KaiTi}
\setCJKfamilyfont{li}{LiSu}
\fi
\newcommand{\song}{\CJKfamily{song}}    % 宋体   (Windows自带simsun.ttf)
\newcommand{\fs}{\CJKfamily{fs}}        % 仿宋体 (Windows自带simfs.ttf)
\newcommand{\kai}{\CJKfamily{nkai}}      % 楷体   (Windows自带simkai.ttf)
\newcommand{\hei}{\CJKfamily{hei}}      % 黑体   (Windows自带simhei.ttf)
\newcommand{\li}{\CJKfamily{li}}        % 隶书   (Windows自带simli.ttf)

%%%%%%%%%%%%%%%%%%%%%%%%%%%%%%%%%%%%%%%%%%%%%%%%%%%%%%%%%%%
% 重定义字号命令
%%%%%%%%%%%%%%%%%%%%%%%%%%%%%%%%%%%%%%%%%%%%%%%%%%%%%%%%%%%
\newcommand{\yihao}{\fontsize{26pt}{26pt}\selectfont}       % 一号, 1.倍行距
\newcommand{\erhao}{\fontsize{22pt}{22pt}\selectfont}       % 二号, 1.倍行距
\newcommand{\xiaoer}{\fontsize{18pt}{18pt}\selectfont}      % 小二, 单倍行距
\newcommand{\sanhao}{\fontsize{16pt}{16pt}\selectfont}      % 三号, 1.倍行距
\newcommand{\xiaosan}{\fontsize{15pt}{15pt}\selectfont}     % 小三, 1.倍行距
\newcommand{\sihao}{\fontsize{14pt}{14pt}\selectfont}       % 四号, 1.0倍行距
\newcommand{\banxiaosi}{\fontsize{13pt}{13pt}\selectfont}   % 半小四, 1.0倍行距
\newcommand{\xiaosi}{\fontsize{12pt}{12pt}\selectfont}      % 小四, 1.倍行距
\newcommand{\dawuhao}{\fontsize{11.5pt}{11.5pt}\selectfont} % 大五号, 单倍行距
\newcommand{\wuhao}{\fontsize{10.5pt}{10.5pt}\selectfont}   % 五号, 单倍行距
\newcommand{\xiaowu}{\fontsize{9.5pt}{9.5pt}\selectfont}    % 五号, 单倍行距
\newcommand{\banbanxiaosi}{\fontsize{12pt}{12pt}\selectfont}% 半半小四, 1.0倍行距


%避免宏包 hyperref 和 arydshln 不兼容带来的目录链接失效的问题。
\def\temp{\relax}
\let\temp\addcontentsline
\gdef\addcontentsline{\phantomsection\temp}
\newcommand*{\subfigencaptionlist}{} % 子图形加入目录时用

\makeatletter
\gdef\hitempty{}

\newcommand{\mr}[1]{\mathrm{#1}} %定义新命令，用\mr来代替\mathrm
\def \ReferenceEName {References} %%定义参考文献的标题
\def \ReferenceCName {参考文献}

%定义图表章节双标题命令
\newcommand{\figenname}{Fig.}
\newcommand{\listfigenname}{List of Figures}
\newfloatlist[chapter]{figen}{fen}{\listfigenname}{\figenname}
\newfixedcaption{\figencaption}{figen}
\renewcommand{\thefigen}{\thechapter.\arabic{figure}}
\renewcommand{\@cftmakefentitle}{\chapter*{\listfigenname\@mkboth{\bfseries\listfigenname}{\bfseries\listfigenname}}}


\newcommand{\FigureBiCaption}[2]
{\renewcommand{\figurename}{图}
\caption{\protect\setlength{\baselineskip}{1.5em}#1} %\protect\setlength{\baselinestretch}{1.3}\selectfont
\vspace{-1.3ex}%-0.5ex
\figencaption{\protect\setlength{\baselineskip}{1.5em}#2}%
\vspace{-3.4mm}
%%
%%其子图形加入目录
\def\hittemp{}
 \@for \hittemp:=\subfigencaptionlist \do {%
        \ifx \hitempty\hittemp\relax \else
          \addcontentsline{fen}{subfigen}{\protect\numberline\hittemp}
        \fi}
 \gdef\subfigencaptionlist{}
}


\setcounter{fendepth}{2} %英文图形目录的深度 1(只有一级目录) 2(有两级目录)
\setcounter{lofdepth}{2} %中文图形目录的深度 1(只有一级目录) 2(有两级目录)
\renewcommand*{\l@subfigure}{\@dottedxxxline{\ext@subfigure}{2}{3.8em}{1.5em}} %中文图形目录 subfigure
\gdef\l@subfigen{\@dottedtocline{2}{3.8em}{1.5em}}%英文图形目录 latex
\newif\ifsubfigtoc
\ifnum \tw@ > \@nameuse{c@fendepth} \subfigtocfalse \else \subfigtoctrue \fi
\newbox\tempbox
\newcommand{\SubfigEnCaption}[1]
{\ifsubfigtoc
    %加入目录这个动作，一定要在 父图 之后,所在先暂存在 subfigencaptionlist
    \protected@xdef\subfigencaptionlist{\subfigencaptionlist,%
        {{\thesubfigure}\protect\ignorespaces{#1}}}
\else
    \relax
\fi
%产生caption
\vspace{1pt}
\sbox{\tempbox}{\thesubfigure\hskip\subfiglabelskip #1}%
\ifthenelse{\lengthtest{\wd\tempbox > \linewidth}}%
{\\[-20pt]\parbox[t]{\linewidth}{\flushleft\noindent\wuhao\selectfont\thesubfigure\hskip\subfiglabelskip \centering#1\hangafter=1\hangindent=15pt}}%
{\\\centerline{\wuhao\selectfont\thesubfigure\hskip\subfiglabelskip #1}}
}

%\newcommand{\SubfigureCaption}[2]  % Two Parameters, the first one is the width of the subfigure,
%{
%\addtocounter{subfigure}{-1}       % the second one is the caption of the subfigure
%\vspace{-2ex}
%\subfigure[#2]{\rule{#1}{0pt}}
%}

\newcommand{\tblenname}{Table} %define tbl instead of table
\newcommand{\listtblenname}{List of Tables}
\newfloatlist[chapter]{tblen}{ten}{\listtblenname}{\tblenname}
\newfixedcaption{\tblencaption}{tblen}
\renewcommand{\thetblen}{\thechapter.\arabic{table}}% 将tblen换成table，因为table和tablen编号一致，而tablen在\longbitoccaption定义中无效。
\renewcommand{\@cftmaketentitle}{\chapter*{\listtblenname\@mkboth{\bfseries\listtblenname}{\bfseries\listtblenname}}}


\newcommand{\TableBiCaption}[2]
{
\renewcommand{\tablename}{表}
\caption{\protect\setlength{\baselineskip}{1.5em}#1}
\vspace{-2ex}
\tblencaption{\protect\setlength{\baselineskip}{1.5em}#2}
\vspace{1ex}
}

%%%% 长表格的caption在中英文表格目录中正常显示
\def\@cont@LT@LTBiToeCaption#1[#2]#3{%
  \LT@makecaption#1\fnum@table{#3}%
  \def\@tempa{#2}%
  \ifx\@tempa\@empty\else
    {\let\\\space
      %\phantomsection
      \addcontentsline{ten}{tblen}{\protect\numberline{\thetable}{#2}}}%%\addcontentsline{lot}{table}{\protect\numberline{}{#2}}}%
  \fi}
\def\LT@c@ption#1[#2]#3{%
  \LT@makecaption#1\fnum@table{#3}%
  \def\@tempa{#2}%
  \ifx\@tempa\@empty\else
     {\let\\\space
     %\phantomsection
     \addcontentsline{lot}{table}{\protect\numberline{\thetable}{#2}}}%
  \fi}
\let\@cont@oldLT@c@ption\LT@c@ption
\newcommand*{\LTBiTocCaption}[5]{
  \@if@contemptyarg{#1}{\caption{#2}}{\caption[#1]{#2}}%
  \global\let\@cont@oldtablename\tablename
  \gdef\tablename{Table} %#3
  \global\let\LT@c@ption\@cont@LT@LTBiToeCaption
  \\
  \@if@contemptyarg{#4}{\caption{#5}}{\caption[#4]{#5}}%
  \global\let\tablename\@cont@oldtablename
  \global\let\LT@c@ption\@cont@oldLT@c@ption}

\renewcommand{\cfttblendotsep}{1} %自定义图表目录中的点间距大小
\renewcommand{\cftfigendotsep}{1}

%\renewcommand{\tablename}{表}  %jdg提供的一种方法，英文长表会添加到中文目录中去，上面定义
%\newcommand{\LTBiCaption}[2]   %\bicaptiontwotoc 主要就是解决这个问题。
%{%
%\caption{#1} \gdef\tablename{Table}
%\\ %[-3.5ex]
%\caption{#2}
%\gdef\tablename{表}\\ %[-1.5ex]
%}

%%%---公式中符号描述----start----
%\begin{formulasymb}{式中}{-3pt}%-3pt,-20pt调与上方的间距。
%  \fdesfirst{第一标签}{控制控制控制控制控制}
%  \fdes{其他标签}{控制控制控制控制控制}
%\end{formulasymb}
\newenvironment{formulades}[1]%
{\noindent\begin{list}{}{%
\setlength\topsep{0pt}
\settowidth{\labelwidth}{#1}
\setlength{\labelsep}{1mm}
\setlength{\leftmargin}{\labelwidth+\labelsep}
}}{\end{list}}
\newenvironment{formulasymb}[2]%-\!-\!-\!-
{\vspace*{#2}\newcommand{\fdesfirst}[2]%
{\begin{formulades}{#1\hspace*{26pt}##1~\cdash}\item[#1\hspace*{26pt}##1~\cdash]{##2}\end{formulades}\vspace*{-21pt}}%自己调距
\newcommand{\fdes}[2]{\begin{formulades}{#1\hspace{26pt}##1~\cdash}\item[##1~\cdash]{##2}\end{formulades}\vspace*{-21pt}}}%自己调距
{\vspace{21pt}\relax}%21pt调距
%%----公式中符号描述----end-----

%% ---- 左对齐的公式 start-----   \begin{flualign} a=c \end{flualign}
\newenvironment{flualign}{%
    \@fleqntrue
    \@mathmargin = -1sp
    \@mathmargin\leftmargini minus\leftmargini
    \let\mathindent=\@mathmargin
  \start@align\@ne\st@rredfalse\m@ne
}{%
  \math@cr \black@\totwidth@
  \egroup
  \ifingather@
    \restorealignstate@
    \egroup
    \nonumber
    \ifnum0=`{\fi\iffalse}\fi
  \else
    $$%
  \fi
  \ignorespacesafterend
  \@fleqnfalse
}
%%  ---- 左对齐的公式  end----

%重新定义BiChapter命令，可实现标题手动换行，但不影响目录
\def\BiChapter{\relax\@ifnextchar [{\@BiChapter}{\@@BiChapter}}
\def\@BiChapter[#1]#2#3{\chapter[#1]{#2}
    \addcontentsline{toe}{chapter}{\bfseries Chapter \thechapter\hspace{0.5em} #3}}
\def\@@BiChapter#1#2{\chapter{#1}
    \addcontentsline{toe}{chapter}{\bfseries Chapter \thechapter\hspace{0.5em}{\boldmath #2}}}
%\newcommand{\BiChapter}[2]
%{
%    \chapter{#1}
%    \addcontentsline{toe}{chapter}{\bfseries Chapter \thechapter\hspace{0.5em} #2}
%}

\newcommand{\BiSection}[2]
{   \section{#1}
    \addcontentsline{toe}{section}{\protect\numberline{\csname thesection\endcsname}#2}
}

\newcommand{\BiSubsection}[2]
{    \subsection{#1}
    \addcontentsline{toe}{subsection}{\protect\numberline{\csname thesubsection\endcsname}#2}
}

\newcommand{\BiSubsubsection}[2]
{    \subsubsection{#1}
    \addcontentsline{toe}{subsubsection}{\protect\numberline{\csname thesubsubsection\endcsname}#2}
}

\newcommand{\BiAppendixChapter}[2] % 该附录命令适用于发表文章，简历等
{\phantomsection
\markboth{#1}{#1}%\markboth{\MakeUppercase{#1}}{\MakeUppercase{#1}}
\addcontentsline{toc}{chapter}{\hei #1}
\addcontentsline{toe}{chapter}{\bfseries #2}  \chapter*{#1}
}

\newcommand{\BiAppChapter}[2]    % 该附录命令适用于有章节的完整附录
{\phantomsection  \chapter{#1}   %\markboth{\MakeUppercase{#1}}{\MakeUppercase{#1}} %为了winedt中project tree中toc正确显示，不要挪到下一行；
%\addcontentsline{toc}{chapter}{\hei #1}
 \addcontentsline{toe}{chapter}{\bfseries Appendix \thechapter~~#2}
}

\renewcommand{\thefigure}{\arabic{chapter}.\arabic{figure}}%使图编号为 7.1 的格式 %\protect{~}
%\renewcommand\fnum@figure{\figurename\nobreakspace\thefigure\protect{~~~~~~~~~}} %

\renewcommand{\thesubfigure}{\alph{subfigure})}%使子图编号为 a)的格式
\renewcommand{\p@subfigure}{\thefigure(} %%使子图引用为 7.1(a) 的格式
%\renewcommand{\thesubfigure}{\alph{subfigure}}
%\renewcommand{\p@subfigure}{\thefigure} %%使子图引用为 7.1a 的格式
%\renewcommand{\@thesubfigure}{\thesubfigure)\hskip\subfiglabelskip}%使子图编号为 a)的格式

\renewcommand{\thetable}{\arabic{chapter}.\arabic{table}}%%使表编号为 7.1 的格式
\renewcommand{\theequation}{\arabic{chapter}.\arabic{equation}}%%使公式编号为 7.1 的格式

% 设置算法标题形式，由“Algorithm 2.1：算法标题” 改为“算法 2.1 算法标题”
\renewcommand{\algorithmcfname}{算法}
\setlength\AlCapSkip{1.2ex}
\SetAlgoSkip{1pt}
\renewcommand{\algocf@captiontext}[2]{#1\algocf@typo ~ \AlCapFnt{}#2} % text of caption
\expandafter\ifx\csname algocf@within\endcsname\relax% if \algocf@within doesn't exist
\renewcommand\thealgocf{\@arabic\c@algocf} % and the way it is printed
\else%                                    else
\renewcommand\thealgocf{\csname the\algocf@within\endcsname.\@arabic\c@algocf}
\fi

\makeatother
%定义 学科 学位
\def \xuekeEngineering {Engineering}
\def \xuekeScience {Science}
\def \xuekeManagement {Management}
\def \xuekeArts {Arts}

\ifx \xueke \xuekeEngineering
\newcommand{\cxueke}{工学}
\newcommand{\exueke}{Engineering}
\fi

\ifx \xueke \xuekeScience
\newcommand{\cxueke}{理学}
\newcommand{\exueke}{Science}
\fi

\ifx \xueke \xuekeManagement
\newcommand{\cxueke}{管理学}
\newcommand{\exueke}{Management}
\fi

\ifx \xueke \xuekeArts
\newcommand{\cxueke}{文学}
\newcommand{\exueke}{Arts}
\fi

\newcommand{\cdash}{\mbox{—\!\!\!\!—\!\!\!\!—}}%输入中文破折号的命令
\newcommand{\dif}{\mathrm{d}}%在数学模式中输入微分dx

%将参考文献在文中的引用改成上标形式 --------by tina
%\newcommand{\upcite}[1]{\textsuperscript{\textsuperscript{\cite{#1}}}}
\newcommand{\upcite}[1]{\textsuperscript{\cite{#1}}}  % 文本格式定义
% !Mode:: "TeX:UTF-8"
% 若xelatex编译非UTF8文件，需在每个文件中指定字符编码;
% main.tex中手动制定了\atemp和\usewhat参数；
\ifx\atempxetex\usewhat
%\XeTeXinputencoding "gbk"
\fi

%%%%%%%%%%%%%%%%%%%%%%%%%%%%%%%%%%%%%%%%%%%%%%%%%%%%%%%%%%%
%允许公式换页显示,否则大型推导公式都在一页内，
%一页显示不下放到第二页，导致很大的空白空间，很不好看
\allowdisplaybreaks[4]

%%%%%%%%%%%%%%%%%%%%%%%%%%%%%%%%%%%%%%%%%%%%%%%%%%%%%%%%%%%
%下面这组命令使浮动对象的缺省值稍微宽松一点，从而防止幅度
%对象占据过多的文本页面，也可以防止在很大空白的浮动页上放置很小的图形。
\renewcommand{\topfraction}{0.9999999}
\renewcommand{\textfraction}{0.0000001}
\renewcommand{\floatpagefraction}{0.9999}

%%%%%%%%%%%%%%%%%%%%%%%%%%%%%%%%%%%%%%%%%%%%%%%%%%%%%%%%%%%
% 重定义一些正文相关标题
%%%%%%%%%%%%%%%%%%%%%%%%%%%%%%%%%%%%%%%%%%%%%%%%%%%%%%%%%%%
\theoremstyle{plain} \theorembodyfont{\song\rmfamily}
\theoremheaderfont{\hei\rmfamily} %\theoremseparator{:}
\newtheorem{definition}{\hei 定义}[chapter]
\newtheorem{example}{\hei 例}[chapter]
\newtheorem{algo}{\hei 算法}[chapter]
\newtheorem{theorem}{\hei 定理}[chapter]
\newtheorem{axiom}{\hei 公理}[chapter]
\newtheorem{proposition}{\hei 命题}[chapter]
\newtheorem{lemma}{\hei 引理}[chapter]
\newtheorem{corollary}{\hei 推论}[chapter]
\newtheorem{remark}{\hei 注解}[chapter]
%\newtheorem{proposition}[definition]{\hei 命题}
%\newtheorem{lemma}[definition]{\hei 引理}
%\newtheorem{exercise}[definition]{}
%\newtheorem{corollary}[definition]{\hei 推论}
%\newtheorem{remark}[definition]{\hei 注解}
%%%%%%%%%%%%%%%%%%%%%%%%%%%%%%%%%%%%%%%%%%%%%%%%%%%%%%%%%%%%%%%%%%%
%解决原proof定理环境的两个问题：
%  1. proof 中的item缩进不对
%  2. proof 中的最后一个公式下出现一个黑方块。
%\theoremsymbol{$\blacksquare$}
%\newtheorem{proof}{\hei 证明}
\newenvironment{proof}{\noindent{\hei 证明：}}{\hfill $ \square $ \vskip 4mm}
\theoremsymbol{$\square$}


%%%%%%%%%%%%%%%%%%%%%%%%%%%%%%%%%%%%%%%%%%%%%%%%%%%%%%%%%%%
% 用于中文段落缩进 和正文版式
%%%%%%%%%%%%%%%%%%%%%%%%%%%%%%%%%%%%%%%%%%%%%%%%%%%%%%%%%%%

\ifx\atempxetex\usewhat
\newcommand{\CJKcaption}[1]{
  \ifx\CJK@actualBinding \@empty
    \PackageError{CJK}{
      You must be inside of a CJK environment to use \protect\CJKcaption}{}
  \else
    \makeatletter
    \InputIfFileExists{#1.cpx}{}{
      \PackageError{CJK}{
        Can't find #1.cpx}{
        The default captions are used if you continue.}}
    \makeatother
  \fi}
\CJKcaption{gb_452}
\else
\CJKcaption{gb_452}
\newlength \CJKtwospaces
\def\CJKindent{
    \settowidth\CJKtwospaces{\CJKchar{"0A1}{"0A1}\CJKchar{"0A1}{"0A1}}%
    \parindent\CJKtwospaces
}
\CJKtilde  \CJKindent
\fi
\renewcommand\CJKprechaptername{}
\renewcommand\CJKchaptername{}
\setlength{\parindent}{26pt} %由于工大论文的每行的字距加大了，需要增加段首缩2pt

\renewcommand\contentsname{\hei 目~~~~录}

%%%%%%章节标题为“1”的形式
\renewcommand\chaptername{\arabic{chapter}}
%%%%%%%%%%%%%%%%%%%%%%%%%%%%%%%%%%%%%%%%%%%%%%%%%%
%定义段落章节的标题和目录项的格式
%%%%%%%%%%%%%%%%%%%%%%%%%%%%%%%%%%%%%%%%%%%%%%%%%%
\setcounter{secnumdepth}{4} \setcounter{tocdepth}{2}

\titleformat{\chapter}[hang]{\xiaosan\bf\filcenter\song\sf\boldmath}{\xiaosan\chaptertitlename}{1em}{\xiaosan}
\titlespacing{\chapter}{0pt}{18pt}{18pt}

\titleformat{\section}[hang]{\song\sihao\sf\bf\boldmath}{\song\sihao\sf\bf\thesection}{0.5em}{}
\titlespacing{\section}{0pt}{12pt}{12pt}

\titleformat{\subsection}[hang]{\song\sf\xiaosi\bf\boldmath}{\xiaosi\thesubsection}{0.5em}{}
\titlespacing{\subsection}{0pt}{6pt}{6pt}

\titleformat{\subsubsection}[runin]{\song\sf\xiaosi\boldmath}{\xiaosi\thesubsubsection}{0.5em}{}[\;\;]
\titlespacing{\subsubsection}{0pt}{3pt}{2pt}

% 缩小目录中各级标题之间的缩进，使它们相隔一个字符距离，也就是12pt
\makeatletter
\renewcommand*\l@section{\@dottedtocline{1}{12pt}{18pt}}
\renewcommand*\l@subsection{\@dottedtocline{2}{24pt}{27pt}}
\renewcommand*\l@subsubsection{\@dottedtocline{3}{36pt}{39pt}}
\renewcommand*\l@paragraph{\@dottedtocline{4}{48pt}{48pt}}
\renewcommand*\l@subparagraph{\@dottedtocline{5}{60pt}{60pt}}

%控制中文目录
\dottedcontents{chapter}[3.4em]{\vspace{0.5em}\hspace{-3.4em}\hei \bf\boldmath}{0.0em}{5pt}% 章标题后用粗点
%\titlecontents{chapter}[3.92em]{\vspace{0.5em}\hspace{-3.92em}\hei \bf\boldmath}{\contentslabel{0em}}{\hspace*{-0em}}{\normalfont\titlerule*[5pt]{.}\contentspage} %章标题后用细点
\dottedcontents{section}[1.6cm]{}{1.8em}{5pt}
\dottedcontents{subsection}[2.840cm]{}{2.7em}{5pt}
\dottedcontents{subsubsection}[3.78cm]{}{3.4em}{5pt}

%%%%%%%%%%%%%%%%%%%%%%%%%%%%%%%%%%%%%%%%%%%%%%%%%%%%%%%
% 定义页眉和页脚 使用fancyhdr 宏包
%%%%%%%%%%%%%%%%%%%%%%%%%%%%%%%%%%%%%%%%%%%%%%%%%%%%%%%%
\newcommand{\makeheadrule}{%
\makebox[-3pt][l]{\rule[.7\baselineskip]{\headwidth}{0.4pt}}
\rule[0.85\baselineskip]{\headwidth}{0pt}\vskip-.8\baselineskip}%控制页眉上的线粗细
\renewcommand{\headrule}{%
    {\if@fancyplain\let\headrulewidth\plainheadrulewidth\fi
     \makeheadrule}}

\pagestyle{fancyplain}

%去掉章节标题中的数字
%%不要注销这一行，否则页眉会变成：“第1章1  绪论”样式
\renewcommand{\chaptermark}[1]{\markboth{\chaptertitlename~~ \ #1}{}}
 \fancyhf{}

%在book文件类别下,\leftmark自动存录各章之章名,\rightmark记录节标题
%% 页眉字号 南理工要求小五宋体
%根据单双面打印设置不同的页眉；
\ifoneortwoside
  \fancyhead{} % clear all fields
  %\fancyhead[LO]{\CJKfamily{song}\xiaowu{} \cxueke\cxuewei 学位论文}
  \fancyhead[LO]{\CJKfamily{song}\xiaowu{} \cxuewei 学位论文}
  \fancyhead[RO]{\CJKfamily{song}\xiaowu XXXXXXXXXXXXXXXXXXx}
  \fancyhead[LE]{\CJKfamily{song}\xiaowu \leftmark}%
  %\fancyhead[RE]{\CJKfamily{song}\xiaowu  \cxueke\cxuewei 学位论文}%
  \fancyhead[RE]{\CJKfamily{song}\xiaowu{} \cxuewei 学位论文}
  \fancyfoot[RO,LE]{\xiaowu ~\thepage~} %页码左右两边有一小段横线%\if@mainmatter \fi
\else
\fancyhead{} % clear all fields
  %\fancyhead[LO]{\CJKfamily{song}\xiaowu{} \cxueke\cxuewei 学位论文}
  \fancyhead[LO]{\CJKfamily{song}\xiaowu{} \cxuewei 学位论文}
  \fancyhead[RO]{\CJKfamily{song}\xiaowu{} 基于多核学习的聚类与层次分类算法研究及应用}
  \fancyhead[LE]{\CJKfamily{song}\xiaowu\leftmark }%
  %\fancyhead[RE]{\CJKfamily{song}\xiaowu  \cxueke\cxuewei 学位论文 }%
  \fancyhead[RE]{\CJKfamily{song}\xiaowu{} \cxuewei 学位论文}
  \fancyfoot[RO,LE]{\xiaowu ~\thepage~} %\if@mainmatter \fi
\fi

\renewcommand\frontmatter{%
    \cleardoublepage
  \@mainmatterfalse
  \pagenumbering{Roman}}
%%%%%%%%%%%%%%%%%%%%%%%%%%%%%%%%%%%%%%%%%%%%%%%%%%%%%%%%
% 设置行距和段落间垂直距离
%%%%%%%%%%%%%%%%%%%%%%%%%%%%%%%%%%%%%%%%%%%%%%%%%%%%%%%%
\renewcommand{\CJKglue}{\hskip 0.3pt plus 0.08\baselineskip}%加大字间距，使每行33个字
%\setlength{\belowcaptionskip}{10pt}   % 加大标题和表格之间的距离 \abovecaptionskip 默认是10pt
\setlength{\parskip}{3pt plus1pt minus1pt} % 段落之间的竖直距离
\renewcommand{\baselinestretch}{1.6}% 定义行距
%%%%%%%%%%%%%%%%%%%%%%%%%%%%%%%%%%%%%%%%%%%%%%%%%%%%%%%%
% 调整列表环境的垂直间距
%%%%%%%%%%%%%%%%%%%%%%%%%%%%%%%%%%%%%%%%%%%%%%%%%%%%%%%%
\setitemize{itemindent=38pt,leftmargin=0pt,itemsep=-0.4ex,listparindent=26pt,partopsep=0pt,parsep=0.5ex,topsep=-0.25ex}
\setenumerate{itemindent=38pt,leftmargin=0pt,itemsep=-0.4ex,listparindent=26pt,partopsep=0pt,parsep=0.5ex,topsep=-0.25ex}
\setdescription{itemindent=38pt,leftmargin=0pt,itemsep=-0.4ex,listparindent=26pt,partopsep=0pt,parsep=0.5ex,topsep=-0.25ex}

\newcommand{\ucite}[1]{$^{\mbox{\scriptsize \cite{#1}}}$} % 增加 \ucite 命令使显示的引用为上标形式
\newcommand{\citeup}[1]{$^{\mbox{\scriptsize \cite{#1}}}$} % for WinEdt users

%%%%%%%%%%%%%%%%%%%%%%%%%%%%%%%%%%%%%%%%%%%%%%%%%%%%%%%%%%%
% 定制浮动图形和表格标题样式 %这里用ccaption完全代替了caption2的功能
\captionstyle{\centering}   %不同的图标题形式采用不同的命令
%\indentcaption{0pt}           %参见ccaption
\hangcaption
\captionnamefont{\CJKfamily{song}\rmfamily\wuhao\selectfont}
\captiontitlefont{\CJKfamily{song}\rmfamily\wuhao\selectfont}
\captiondelim{~} %~

%%%%%%%%%%%%%%%%%%%%%%%%%%%%%%%%%%%%%%%%%%%%%%%%%%%%%%%
% 定义题头格言的格式
% 用法 \begin{Aphorism}{author}
%         aphorism
%      \end{Aphorism}
\newsavebox{\AphorismAuthor}
\newenvironment{Aphorism}[1]
{\vspace{0.5cm}\begin{sloppypar} \slshape
\sbox{\AphorismAuthor}{#1}
\begin{quote}\small\itshape }
{\\ \hspace*{\fill}------\hspace{0.2cm} \usebox{\AphorismAuthor}
\end{quote}
\end{sloppypar}\vspace{0.5cm}}

%自定义一个空命令，用于注释掉文本中不需要的部分。
\newcommand{\comment}[1]{}

\renewcommand\contentsname{\hei 目~~~~录}
\renewcommand\listfigurename{\hei 插图目录}
\renewcommand\listtablename{\hei 表格目录}

%%%%%%将章标题中的中文数字（一、二、三）变为阿拉伯数字(1,2,3)
\renewcommand\CJKthechapter{%\CJKnumber
{\@arabic\c@chapter}}

%%%%%%不要拉大行距使得页面充满
\raggedbottom

% This is the flag for longer version
\newcommand{\longer}[2]{#1}

\newcommand{\ds}{\displaystyle}

% define graph scale
\def\gs{1.0}

%%%%%%%%%%%%%%%%%%%%%%%%%%%%%%%%%%%%%%%%%%%%%%%%%%%%%%%%%%%%%%%%%%%%%%
% 自定义项目列表标签及格式 \begin{hitlist} 列表项 \end{hitlist}
%%%%%%%%%%%%%%%%%%%%%%%%%%%%%%%%%%%%%%%%%%%%%%%%%%%%%%%%%%%%%%%%%%%%%%
\newcounter{hitctr} %自定义新计数器
\newenvironment{hitlist}{%%%%%定义新环境
\begin{list}{{\hei (\arabic{hitctr})}} %%标签格式
    {
     \usecounter{hitctr}
     \setlength{\leftmargin}{0cm}     %左边界
     \setlength{\parsep}{0ex}         %段落间距
     \setlength{\topsep}{0pt}         %列表到上下文的垂直距离
     \setlength{\itemsep}{0ex}        %标签间距
     \setlength{\labelsep}{0.3em}     %标号和列表项之间的距离,默认0.5em
     \setlength{\itemindent}{46pt}    %标签缩进量
     \setlength{\listparindent}{27pt} %段落缩进量
    }}
{\end{list}}%%%%%

%%%%%%%%%%%%%%%%%%%%%%%%%%%%%%%%%%%%%%%%%%%%%%%%%%%%%%%%%%%%%%%%%%%%%%
% 自定义项目列表标签及格式 \begin{publist} 列表项 \end{publist}
%%%%%%%%%%%%%%%%%%%%%%%%%%%%%%%%%%%%%%%%%%%%%%%%%%%%%%%%%%%%%%%%%%%%%%
\newcounter{pubctr} %自定义新计数器
\newenvironment{publist}{%%%%%定义新环境
\begin{list}{\arabic{pubctr}} %%标签格式
    {
     \usecounter{pubctr}
     \setlength{\leftmargin}{2em}     % 左边界 \leftmargin =\itemindent + \labelwidth + \labelsep
     \setlength{\itemindent}{0em}     % 标号缩进量
     \setlength{\labelwidth}{1em}     % 标号宽度
     \setlength{\labelsep}{1em}       % 标号和列表项之间的距离,默认0.5em
     \setlength{\rightmargin}{0em}    % 右边界
     \setlength{\topsep}{0ex}         % 列表到上下文的垂直距离
%     \setlength{\partopsep}{0ex}      % 列表是一个新的段落时，加的额外到上下文的距离
     \setlength{\parsep}{0ex}         % 段落间距
     \setlength{\itemsep}{0ex}        % 标签间距
     \setlength{\listparindent}{26pt} % 段落缩进量
    }}
{\end{list}}%%%%%

%%%%%%%%%%%%%%%%%%%%%%%%%%%%%%%%%%%%%%%%%%%%%%%%%%%%%%%%%%%%%%%%%%%%%%
% 自定义项目列表标签及格式 \begin{gagalist1} 列表项 \end{gagalist1}
%%%%%%%%%%%%%%%%%%%%%%%%%%%%%%%%%%%%%%%%%%%%%%%%%%%%%%%%%%%%%%%%%%%%%%
\newcounter{gagactr1} %自定义新计数器
\newenvironment{gagalist1}{%%%%%定义新环境
\begin{list}{\hei(\arabic{gagactr1})} %%标签格式
    {
     \usecounter{gagactr1}
     \setlength{\leftmargin}{2em}     % 列表与左边距之间的水平距离，值为非负
     \setlength{\itemindent}{0em}     % 每一 item 第一行的缩进
     \setlength{\labelwidth}{1em}     % 标号宽度
     \setlength{\labelsep}{1em}       % 标号和列表项之间的距离,默认0.5em
     \setlength{\rightmargin}{0em}    % 右边界
     \setlength{\topsep}{0ex}         % 列表到上下文的垂直距离
%     \setlength{\partopsep}{0ex}      % 列表是一个新的段落时，加的额外到上下文的距离
     \setlength{\parsep}{0ex}         % 段落间距
     \setlength{\itemsep}{0ex}        % 标签间距
     \setlength{\listparindent}{26pt} % 每一 item 第一行之后各行的缩进
    }}
{\end{list}}%%%%%

%%%%%%%%%%%%%%%%%%%%%%%%%%%%%%%%%%%%%%%%%%%%%%%%%%%%%%%%%%%%%%%%%%%%%%
% 自定义项目列表标签及格式 \begin{gagalist1} 列表项 \end{gagalist1}
%%%%%%%%%%%%%%%%%%%%%%%%%%%%%%%%%%%%%%%%%%%%%%%%%%%%%%%%%%%%%%%%%%%%%%
% \newcounter{gagactr2} %自定义新计数器
\newenvironment{gagalist2}{%%%%%定义新环境
\begin{list}{} %%标签格式
    {
     % \usecounter{gagactr2}
     \setlength{\leftmargin}{1em}     % 列表与左边距之间的水平距离，值为非负
     \setlength{\itemindent}{0em}     % 每一 item 第一行的缩进
     \setlength{\labelwidth}{0.5em}     % 标号宽度
     \setlength{\labelsep}{0.5em}       % 标号和列表项之间的距离,默认0.5em
     \setlength{\rightmargin}{0em}    % 右边界
     \setlength{\topsep}{0ex}         % 列表到上下文的垂直距离
%     \setlength{\partopsep}{0ex}      % 列表是一个新的段落时，加的额外到上下文的距离
     \setlength{\parsep}{0ex}         % 段落间距
     \setlength{\itemsep}{0ex}        % 标签间距
     \setlength{\listparindent}{26pt} % 每一 item 第一行之后各行的缩进
    }}
{\end{list}}%%%%%

%%%%%%%%%%%%%%%%%%%%%%%%%%%%%%%%%%%%%%%%%%%%%%%%%%%%%%%%%%%%%%%%%%%%%%
% 默认字体
\renewcommand\normalsize{%
  \@setfontsize\normalsize{12pt}{14.5pt}   % 将正文字体改为小四，固定行距20磅（近似一页34行）---------- by tina
  \setlength\abovedisplayskip{8pt plus 2pt minus 2pt}
  \setlength\abovedisplayshortskip{7pt plus 2pt minus 2pt}
  \setlength\belowdisplayskip{\abovedisplayskip}
  \setlength\belowdisplayshortskip{\abovedisplayshortskip}
  \setlength\jot{6pt}
  \let\@listi\@listI}
\def\defaultfont{\renewcommand{\baselinestretch}{1.37}\normalsize\selectfont}
\predisplaypenalty=0  %公式之前可以换页，公式出现在页面顶部
%%%%%%%%%%%%%%%%%%%%%%%%%%%%%%%%%%%%%%%%%%%%%%%%%%%%%%%%%%
% 定义可以指定宽度的下划线
\def\NJUSTunderline[#1]#2{
\underline{\hbox to #1{\hfill#2\hfill}}}\def\@NJUSTunderline{}
%%%%%%%%%%%%%%%%%%%%%%%%%%%%%%%%%%%%%%%%%%%%%%%%%%%%%%%%%%%%%%%%%%%%%%
% 封面、摘要、版权、致谢格式定义
%--------定义标签对应选项的值-----------------------------
\def\natclassifiedindex#1{\def\@natclassifiedindex{#1}}\def\@natclassifiedindex{}
\def\secretclassifiedindex#1{\def\@secretclassifiedindex{#1}}\def\@secretclassifiedindex{}%密级
\def\internatclassifiedindex#1{\def\@internatclassifiedindex{#1}}\def\@internatclassifiedindex{}

\def\cname#1{\def\@cname{#1}}\def\@cname{}
\def\crole#1{\def\@crole{#1}}\def\@crole{}
\def\school#1{\def\@school{#1}}\def\@school{}
\def\ctitle#1{\def\@ctitle{#1}}\def\@ctitle{}
\def\cdegree#1{\def\@cdegree{#1}}\def\@cdegree{}
\def\csubject#1{\def\@csubject{#1}}\def\@csubject{}
\def\cauthor#1{\def\@cauthor{#1}}\def\@cauthor{}
\def\csupervisor#1{\def\@csupervisor{#1}}\def\@csupervisor{}
%\def\cassosupervisor#1{\def\@cassosupervisor{~ & {\hei 副 \hfill 导 \hfill 师：} & #1\\}}\def\@cassosupervisor{}
%\def\ccosupervisor#1{\def\@ccosupervisor{~ & {\hei 联 \hfill 合\hfill 导 \hfill 师：} & #1\\}}\def\@ccosupervisor{}
\def\cdate#1{\def\@cdate{#1}}\def\@cdate{}
\def\ddate#1{\def\@ddate{#1}}\def\@ddate{}%定义答辩日期
\def\chairman#1{\def\@chairman{#1}}\def\@chairman{}
\def\expositor#1{\def\@expositor{#1}}\def\@expositor{}


%--------------定义标签(label)---------------------
\def\lnatclassifiedindex#1{\def\@lnatclassifiedindex{#1}}\def\@lnatclassifiedindex{}
\def\lsecretclassifiedindex#1{\def\@lsecretclassifiedindex{#1}}\def\@lsecretclassifiedindex{}%密级
\def\linternatclassifiedindex#1{\def\@linternatclassifiedindex{#1}}\def\@linternatclassifiedindex{}

\def\lname#1{\def\@lname{#1}}\def\@lname{}
\def\lauthor#1{\def\@lauthor{#1}}\def\@lauthor{}
\def\lschool#1{\def\@lschool{#1}}\def\@lschool{}
\def\ldegree#1{\def\@ldegree{#1}}\def\@ldegree{}
\def\lsubject#1{\def\@lsubject{#1}}\def\@lsubject{}
\def\lsupervisor#1{\def\@lsupervisor{#1}}\def\@lsupervisor{}
%\def\lassosupervisor#1{\def\@lassosupervisor{~ & {\hei 副 \hfill 导 \hfill 师：} & #1\\}}\def\@lassosupervisor{}
%\def\lcosupervisor#1{\def\@lcosupervisor{~ & {\hei 联 \hfill 合\hfill 导 \hfill 师：} & #1\\}}\def\@lcosupervisor{}
\def\ldate#1{\def\@ldate{#1}}\def\@ldate{}
\def\lddate#1{\def\@lddate{#1}}\def\@lddate{}%定义答辩日期
\def\lchairman#1{\def\@lchairman{#1}}\def\@lchairman{}
\def\lexpositor#1{\def\@lexpositor{#1}}\def\@lexpositor{}

\long\def\cabstract#1{\long\def\@cabstract{#1}}\long\def\@cabstract{}%摘要标签
\def\ckeywords#1{\def\@ckeywords{#1}}\def\@ckeywords{}%关键词标签

%--------------定义英文封面---------------------
\def\etitle#1{\def\@etitle{#1}}\def\@etitle{}
\def\edegree#1{\def\@edegree{#1}}\def\@edegree{}
\def\eaffil#1{\def\@eaffil{#1}}\def\@eaffil{}
\def\esubject#1{\def\@esubject{#1}}\def\@esubject{}
\def\eauthor#1{\def\@eauthor{#1}}\def\@eauthor{}
\def\esupervisor#1{\def\@esupervisor{#1}}\def\@esupervisor{}
%\def\eassosupervisor#1{\def\@eassosupervisor{#1}}\def\@eassosupervisor{}
%\def\eassosupervisor#1{\def\@eassosupervisor{~ & \textbf{Associate Supervisor:} & #1\\}}\def\@eassosupervisor{}
\def\ecosupervisor#1{\def\@ecosupervisor{#1}}\def\@ecosupervisor{}
%\def\ecosupervisor#1{\def\@ecosupervisor{~ & \textbf{Co Supervisor:} & #1\\}}\def\@ecosupervisor{}
\def\edate#1{\def\@edate{#1}}\def\@edate{}
\long\def\eabstract#1{\long\def\@eabstract{#1}}\long\def\@eabstract{}
\long\def\NotationList#1{\long\def\@NotationList{#1}}\long\def\@NotationList{}
\def\ekeywords#1{\def\@ekeywords{#1}}\def\@ekeywords{}
%%%%%%%%%%%%%%%%%%%%%%%%%%%%%%%%%%%%%%%%%%%%%%%%%%%%%%%%%%%%%%%
% 定义封面
\def\makecover{
    \normalbiao %表格字号设置
     %封面一
    \newpage
    \thispagestyle{empty}
    \begin{center}

      \parbox[t][0.6cm][t]{\textwidth}{
      \begin{center} \end{center}}

      \parbox[t][2.2cm][t]{\textwidth}{%
       \song \wuhao
       \@lnatclassifiedindex~
       \NJUSTunderline[100pt]{\@natclassifiedindex }
        \hfill
       \@lsecretclassifiedindex~
       \NJUSTunderline[100pt]{\@secretclassifiedindex}
        \vskip 10pt
        \@linternatclassifiedindex~
        \NJUSTunderline[108pt]{\@internatclassifiedindex}\hfill
       }%

      \parbox[t][2.7cm][b]{\textwidth}{\song\yihao
      \begin{center} {\textbf{学\qquad 位\qquad 论\qquad 文}}\end{center} }

      \setlength{\baselineskip}{1.5\baselineskip}
      \parbox[t][3.0cm][b]{\textwidth}{\erhao
      \begin{center} {\hei\textbf{\@school\@cdegree\@ctitle}}\end{center} }

      \begin{center} {\xiaoer\kai{\NJUSTunderline[100pt]{\textbf{\@cauthor}}}} \end{center}%还可设置上下间隔

      \parbox[t][0.8cm][t]{\textwidth}{\begin{center}  \end{center} }

      \parbox[t][6cm][c]{\textwidth}{
      \begin{center}
        \begin{tabular}{llcll}
        \song\xiaosi\@lsupervisor\@lname~&
        \multicolumn{4}{l}{\NJUSTunderline[320pt]{\kai\sanhao\textbf{\@csupervisor}\hspace*{2em}\kai\sihao\@crole}}\vspace{18pt}\\
        \hspace*{2.5cm}&
        \multicolumn{4}{l}{\NJUSTunderline[320pt]{}}\vspace{18pt}\\
        %\NJUSTunderline[115pt]{\@ccosupervisor}
        %\NJUSTunderline[115pt]{\@cassosupervisor}
        \song\xiaosi\@ldegree &
        \NJUSTunderline[100pt]{\kai\sanhao\textbf{\@cdegree}}&\hspace*{1pt}&
        \song\xiaosi\@lsubject~&
        \NJUSTunderline[100pt]{\kai\sanhao\textbf{\@csubject}}\vspace{18pt}\\
        \song\xiaosi\@ldate~&
        \NJUSTunderline[100pt]{\defaultfont\@cdate}&\hspace*{2pt}&
        \song\xiaosi\@lddate~&
        \NJUSTunderline[100pt]{\defaultfont\@ddate}\vspace{18pt}\\
        \multicolumn{5}{l}{\song\xiaosi\@lschool~~\NJUSTunderline[285pt]{\kai\sanhao\@school}}\vspace{18pt}\\
        \multicolumn{4}{r}{\song\xiaosi\@lchairman}&
        \NJUSTunderline[100pt]{\@chairman}\vspace{18pt}\\
        \multicolumn{4}{r}{\song\xiaosi\@lexpositor}&
        \NJUSTunderline[100pt]{\@expositor}\vspace{18pt}
       \end{tabular}
     \end{center}}

   \end{center}
%%%%%%封面一增加一空白页
  \ifoneortwoside
    \newpage
    ~~~\vspace{1em}
    \thispagestyle{empty}
  \fi
%%%%%%%%%内封
    \begin{titlepage}
    \begin{center}
      \parbox[t][1cm][b]{\textwidth}{\xiaoer
      \begin{center} {\kai \ifxueweimaster\cxueke\fi 中华人民共和国\cxueke\cxuewei 学位论文 }\end{center} }

      \parbox[t][0.8cm][t]{\textwidth}{
      \begin{center} \end{center} }

      \parbox[t][2cm][t]{\textwidth}{\xiaoer
      \begin{center} {\hei\textbf{\@school\@cdegree\@ctitle}}\end{center} }

      \ifxueweidoctor
      \parbox[t][3.8cm][t]{\textwidth}{\xiaoer
      \begin{center} {\hei  \@etitle}\end{center} }
      \else
      \parbox[t][3.8cm][t]{\textwidth}{\centering
        \ }
      \fi

    \parbox[t][2.0cm][t]{\textwidth}{\kai\xiaoer
    \begin{center}
    \begin{tabular}{lcc}
    \textbf{\@lauthor:}&
    \multicolumn{2}{l}{\textbf{\@cauthor}}\\
    \textbf{\@lsupervisor:}&
    \textbf{\@csupervisor }&
    \textbf{\@crole}
    \end{tabular}
    \end{center}}

    \ifxueweidoctor
    %\parbox[t][8.5cm][t]{\textwidth}{\centering
    \parbox[c][8.5cm][c]{\textwidth}{\centering
        \includegraphics[width = 7cm,height=7cm]{xiaohui}}
    \else
    \parbox[t][8.5cm][t]{\textwidth}{\centering \ }
    \fi

    \parbox[t][1.2cm][t]{\textwidth}{\xiaoer
    \begin{center} {\song  \textbf{\@school}}  \end{center} }

    \parbox[t][0.5cm][t]{\textwidth}{
    \begin{center} {\song \xiaoer \textbf{\@cdate}} \end{center} }
   \end{center}

%%%% 内封增加一空白页
    \ifoneortwoside
      \newpage
      ~~~\vspace{1em}
      \thispagestyle{empty}
    \fi

%%%%% 英文封面
    \newpage
    \thispagestyle{empty}
    \begin{center}
    \parbox[t][0.6cm][t]{\textwidth}{
    \begin{center} \end{center}}

   \parbox[t][2.2cm][t]{\textwidth}{}

    \parbox[t][1cm][b]{\textwidth}{\xiaoer
    \begin{center} {PH. D. Dissertation}\end{center} }

    \parbox[t][2mm][t]{\textwidth}{
    \begin{center}  \end{center}}

    \parbox[t][3.0cm][b]{\textwidth}{\erhao
    \begin{center} {\textbf{\@etitle}}\end{center} }

    \parbox[t][3.0cm][t]{\textwidth}{
    \begin{center}  \end{center}}

    \parbox[t][2.5mm][b]{\textwidth}{\xiaoer
    \begin{center} { \emph{By}}\end{center}

    \begin{center} { \emph{\textbf{\@eauthor}}}\end{center}}

    \parbox[t][1.2cm][b]{\textwidth}{\xiaoer
    \begin{center} {\emph{Supervised by Prof.~\textbf{\@esupervisor}}}\end{center}}

    \parbox[t][6.0cm][t]{\textwidth}{
    \begin{center}  \end{center} }

    \parbox[t][2.6cm][b]{\textwidth}{\xiaoer
    \begin{center} {Nanjing University of  Science \& Technology}\end{center}

    \begin{center} {\@edate }\end{center}}

   \end{center}
  \end{titlepage}
%%%%%%英文封面增加一空白页
  \ifoneortwoside
    \newpage
    ~~~\vspace{1em}
    \thispagestyle{empty}
  \fi

 \thispagestyle{empty}
 % !Mode:: "TeX:UTF-8"
% 若xelatex编译非UTF8文件，需在每个文件中指定字符编码;
% main.tex中手动制定了\atemp和\usewhat参数；
\ifx\atempxetex\usewhat
%\XeTeXinputencoding "gbk"
\fi
\defaultfont

\thispagestyle{empty}
%在此节中，由于用了与其它正文不同的字体，此节中用的是
%四号字体，所以要对行距与首行缩进重新定义。
\begin{center}{\sanhao \hei{声\quad 明}}\end{center}

\renewcommand{\baselinestretch}{1.5}\large{}
{\setlength{\parindent}{2em}本学位论文是我在导师的指导下取得的研究成果，尽我所知，在本
 学位论文中，除了加以标注和致谢的部分外，不包含其他人已经发表或
公布过的研究成果，也不包含我为获得任何教育机构的学位或学历而使
用过的材料。与我一同工作的同事对本学位论文做出的贡献均已在论文
中作了明确的说明。}

    \vspace{0.738cm}
    \begin{flushleft}{
    研究生签名：\underline{~~~~~~~~~~~~~~~}~~~~~~~~~~~~~~~~日期:~~~~~~~~~~~年~~~~~月~~~~~日}
    \end{flushleft}

    \vspace{2.214cm}
%%%%%%%%%%%%%%%%%%南京理工大学博（硕）士学位论文使用授权书%%%%%%%%%%%%%%%%%%%
%\phantomsection
    \begin{center}{\sanhao \hei{南京理工大学\cxuewei 学位论文使用授权声明}}
    \end{center}

    \vspace{0.738cm}

\renewcommand{\baselinestretch}{1.5}\large{}
{\setlength{\parindent}{2em}南京理工大学有权保存本学位论文的电子和纸质文档，可以借阅或上网公布本学位论文的部分或全部内容，
可以向有关部门或机构送交并授权其保存、借阅或上网公布本学位论文的部分或全部内容。对于保密论文，
按保密的有关规定和程序处理。}


\vspace{1.476cm}
\begin{flushleft}{
研究生签名：\underline{~~~~~~~~~~~~~~~~~}~~~~~~~~~~~~~~日期:~~~~~~~~~~~年~~~~~月~~~~~日}
\end{flushleft}
\vspace{0.2cm}
\iffalse
\begin{flushleft}{
\hspace*{8pt}导师签名：\underline{~~~~~~~~~~~~~~~~~~~}~~~~~~~~~~~~~~日期:~~~~~~~~~~~年~~~~~月~~~~~日}
\end{flushleft}
\fi
\newpage
   % 原创性申明,即让原创性申明在摘要前显示

%%%%%%增加一空白页
 \ifoneortwoside
    \newpage
    ~~~\vspace{1em}
    \thispagestyle{empty}
  \fi
%%%%%%%%%%%%%%%%%%%   Abstract and keywords  %%%%%%%%%%%%%%%%%%%%%%%修改摘要和关键词???????????
\clearpage \BiAppendixChapter{\textbf{摘\hspace*{2em}要}}{Abstract (in Chinese)} %不要挪到下一行，生成正确的摘要toc
\setcounter{page}{1}
\song \normalsize
\defaultfont
\@cabstract \vspace{1em} \hangafter1\hangindent4.28em\noindent
{\sihao\song \textbf{关键词}} \quad \@ckeywords

%%%%%%%%%%%%%%%%%%%   English Abstract  %%%%%%%%%%%%%%%%%%%%%%%%%%%%%%
\clearpage
\defaultfont \BiAppendixChapter{\textbf{Abstract}}{Abstract (in English)} %不要挪到下一行，生成正确的摘要toc
\@eabstract
\vspace{1em}

\hangafter1\hangindent5.5em\noindent
{\textbf{Keywords}} \quad \@ekeywords
\wuhaobiao  %正文表格设置
}

%%%%%%%%%%%%%%%%%%%%%%%%%%%%%%%%%%%%%%%%%%%%%%%%%%%%%%%%%%%%%%%

\urlstyle{same}  %论文中引用的网址的字体默认与正文中字体不一致，这里修正为一致的。

%主要符号表 \NotationList
\long\def\notation{ \clearpage
\BiAppendixChapter{\textbf{主要符号说明}}{} \normalbiao
\@NotationList \wuhaobiao}

%%% 五号字表格设置 start %%%%%
\gdef\tpltable{\relax}
\let\tpltable\longtable
\gdef\wuhaobiao{%五号字
    \def\tabular{\wuhao\gdef\@halignto{}\@tabular}
    \def\endtabular{\endarray $\egroup \defaultfont}
    \def\longtable{\wuhao\tpltable}
    \def\endlongtable{\adl@LTlastrow \adl@org@endlongtable\defaultfont}
}
\gdef\normalbiao{%正常字号
    \def\tabular{\gdef\@halignto{}\@tabular}
    \def\endtabular{\endarray $\egroup}
    \def\longtable{\tpltable}
    \def\endlongtable{\adl@LTlastrow \adl@org@endlongtable}
}
\wuhaobiao
%%% 五号字表格设置 end %%%%%
%\renewcommand{\arraystretch}{1.4} %表格中行距 ，导致公式 bmatrix 间距增大。

% 表格与下方间距
\renewcommand\endtable{\vspace{-4mm}\end@float}
% 算法与下方间距
\renewcommand\endalgorithm{\@algocf@finish \ifthenelse {\equal {\algocf@float }{figure}}{\end {figure}}{
\@algocf@term@caption \ifthenelse {\boolean {algocf@algoH}}{\end {algocf@Here}}
{\end {algocf}}}\@algocf@term\vspace{-5mm}}

% 插图目录深度设置
\setcounter{lofdepth}{1} %1：只包含Figure的标题，2：包含Figure和Subfigure的标题



\makeatother


%%%%%%%%%%%%%%%%%%%%%%%%%%%%%%%%%%%%%%%%%%%%%%%%%%
% 正文部分
%%%%%%%%%%%%%%%%%%%%%%%%%%%%%%%%%%%%%%%%%%%%%%%%%%
\frontmatter
\sloppy % 解决中英文混排的断行问题，会加入间距，但不会影响断行
%\interlinepenalty -100000

%%% ****** 封面 目录 ******%%%
% !Mode:: "TeX:UTF-8"
% 若xelatex编译非UTF8文件，需在每个文件中指定字符编码;
% main.tex中手动制定了\atemp和\usewhat参数；
\ifx\atempxetex\usewhat
%\XeTeXinputencoding "gbk"
\fi

\newcommand{\chinesethesistitle}{xxxxxxxxxxxxxxxx} %授权书用，无需断行
\newcommand{\englishthesistitle}{xxxxxxxxxxxxxxxxg}
\newcommand{\chinesethesistime}{2013~年~6月}  %封面底部的日期中文形式
\newcommand{\dchinesethesistime}{2013~年~6月} %答辩日期中文形式：2011~年~~月
\newcommand{\englishthesistime}{June, 2013}   %封面底部的日期英文形式
%-----定义标签值-与format.tex文件中的定义相对应----------------
\lnatclassifiedindex{分类号}
\lsecretclassifiedindex{密级}
\linternatclassifiedindex{UDC}

\lsupervisor{指导教师}
\lname{姓名}
\lauthor{作\hspace*{2em}者}
\ldegree{申请学位级别}
\lsubject{专业名称}
\ldate{论文提交日期}
\lddate{论文答辩日期}
\lschool{学位授予单位和日期}
\lchairman{答辩委员会主席}
\lexpositor{评阅人}
\cabstract{摘\hspace*{2em}要}
\ckeywords{关键词}

%----定义标签对应选项的值---------------
\natclassifiedindex{TP309}
\secretclassifiedindex{公开}
\internatclassifiedindex{~681.324} %国际图书分类号

%\school{南京理工大学}
%\cdegree{\cxuewei}
\ctitle{\erhao\hei xxxxxxxxxxxxxxxx}  %封面用论文标题，自己可手动断行
\csubject{xxxxxxxxxxxxxxxx}                 %(~按二级学科填写~)
\cauthor{xxxx}            %作者姓名：XXX
\crole{教授}         %导师职称：教授
\csupervisor{XXXX} %导师名字：XXXX
%\cassosupervisor{某~~~~~~某~~~~教~~授}     %(~如无副导师可以不列此项~)
%\ccosupervisor{某~~某~~某~~~~教~~授~} %(~如无联合培养导师则不列此项~)
\cdate{\chinesethesistime}
\ddate{\dchinesethesistime}%定义答辩日期

\etitle{\englishthesistitle}
\edegree{\exuewei \ of \exueke}
\esubject{Pattern Recognition and Intelligence Systems}  %英文二级学科名
\eaffil{School of Computer Science and Technology}%英文单位 %换行用\newline，不要用\\
\eauthor{\textbf{XXXX}}                   %作者姓名 （英文）：\textbf{Yantingt Lu}
\esupervisor{\textbf{XXXX}}       % 导师姓名 （英文）：\textbf{Jing-Yu Yang}
%\ecosupervisor{Professor X}
%\eassosupervisor{Professor Y}
\edate{\englishthesistime}


\cabstract{
XXXXXXXXXXXXXXXXXXXXXXXXXXXXXXXXXXX

XXXXXXXXXXXXXXXXXXXXXXXXXXXXXXXXXXX

}

%关键词不能放在摘要页的外面
 \ckeywords{XXXXX，\ XXXXX，\ XXXXX， \ XXXXX， \ XXXXX， \ XXXXX， \ XXXXX}



\eabstract{
XXXXXXXXXXXXXXXXXXXXXXXXXXXXXXXXXXX

XXXXXXXXXXXXXXXXXXXXXXXXXXXXXXXXXXX

}

\ekeywords{XXXXX, \ XXXXX, \ XXXXX, \ XXXXX,  \ XXXXX,  \ XXXXX, \ XXXXX}

%主要符号说明
%\NotationList{
%%\vspace*{10pt}%可在主要符号说明与下面内容之间留空
%\begin{center}
%\begin{tabular}{lll}
%符号 & 名字 & 定义\\
%\hline
%$\mathcal {X}$,$\mathcal {Y}$,$\mathcal {Z}$ & 花体大写英文& 张量\\
%$\textbf{A}$, $\textbf{B}$, $\textbf{C}$ & 粗体大写英文& 矩阵 \\
%$\textbf{u}$, $\textbf{v}$, $\textbf{w}$ & 粗体小写英文& 向量 \\
%$l, m, n$ & 常规斜体小写英文& 数量 \\
%$L, M, N$ & 常规斜体大写英文& 数量指标的上界(如$m=1,\ldots,M$) \\
%$\alpha ,\beta ,\theta $ & 小写希腊字母 & 参数 \\
%$I_{n}$ & 专用符号符号 & $n\times n$单位矩阵 \\
%$1_{m\times n}$ & 专用符号 & $m\times n$全一矩阵 \\
%$\epsilon$ & 专用符号 & 误差 \\
%$\lambda $ & 专用符号 & 矩阵的特征值 \\
%$\otimes  $ & 专用符号 & Kronecker积 \\
%$\langle \bullet , \bullet \rangle$ & 专用符号 & 向量内积 \\
%\hline
%\end{tabular}
%\end{center}
%}
\makecover
\clearpage
 % 封面
 \renewcommand{\baselinestretch}{1}
 \fontsize{12pt}{12pt}\selectfont
 \clearpage{\pagestyle{empty}\cleardoublepage}
 \pdfbookmark[0]{目~~~~录}{mulu}
 \tableofcontents    % 中文目录

%%% ****** 图标索引，符号表 ******%%%
 % !Mode:: "TeX:UTF-8"
% 若xelatex编译非UTF8文件，需在每个文件中指定字符编码;
% main.tex中手动制定了\atemp和\usewhat参数；
\ifx\atempxetex\usewhat
%\XeTeXinputencoding "gbk"
\fi

%% 中英文图形表格索引   %硕博士学位论文规范均不要求这一项，请正式打印的时候屏蔽掉这一段；
\ifxueweidoctor
  %\clearpage{\pagestyle{empty}\cleardoublepage}   % 清除目录后面空页的页眉和页脚
\else%
  {\ifoneortwoside\clearpage{\pagestyle{empty}\cleardoublepage}\else\newpage\fi} % 清除目录后面空页的页眉和页脚
\fi
\addcontentsline{toc}{chapter}{\textbf{插图目录}}   % 中文插图加入到中文目录
\listoffigures                                  % 生成中文 图形索引
\ifxueweidoctor                                 %硕士学位论文没有英文目录
%\clearpage{\pagestyle{empty}\cleardoublepage}   % 英文图形索引 右开 ？需要吗？
\fi

\ifxueweidoctor
  %\clearpage{\pagestyle{empty}\cleardoublepage}   % 清除目录后面空页的页眉和页脚
\else%
  {\ifoneortwoside\clearpage{\pagestyle{empty}\cleardoublepage}\else\newpage\fi}   % 清除目录后面空页的页眉和页脚
\fi
\addcontentsline{toc}{chapter}{\textbf{表格目录}}   % 中文表格加入到中文目录
\listoftables                                   % 生成中文 表格索引
\ifxueweidoctor                                 %硕士学位论文没有英文目录
%\clearpage{\pagestyle{empty}\cleardoublepage}   % 英文表格索引 右开 ？需要吗？
\fi
%% 如果不需要图表索引，注释掉上面的即可
  %图表索引, 如果不需要图表索引，注释掉这一句即可；
 %\notation  %主要符号表
 \addtocontents{toc}{\protect\vskip1\baselineskip} % 中文目录增加空行

%%% ****** 清除页眉页脚 ******%%%
\ifxueweidoctor
  \clearpage{\pagestyle{empty}\cleardoublepage}   % 清除目录后面空页的页眉和页脚
\else%
  \ifoneortwoside\clearpage{\pagestyle{empty}\cleardoublepage}\fi  % 清除目录后面空页的页眉和页脚
\fi                                               %  第一章 是否右开

\mainmatter
\defaultfont % 对应于小四的标准字号是 12pt, 可以在正文中用此命令修改所需要字体的的大小

% !Mode:: "TeX:UTF-8"
% 这个是为了WinEdt设置的，它的默认不是UTF8.
% !TeX root = ../main.tex
% 若xelatex编译非UTF8文件，需在每个文件中指定字符编码;
% main.tex中手动制定了\atemp和\usewhat参数；
\ifx\atempxetex\usewhat
%\XeTeXinputencoding "gbk"
\fi

\defaultfont

\titleformat{\chapter}[hang]{\xiaosan\bf\raggedright\hei\sf\boldmath}{\xiaoer\chaptertitlename}{18pt}{\xiaosan}
\titlespacing{\chapter}{0pt}{8pt}{16pt}

\makeatletter
\newskip\@footindent
\@footindent=1em

\renewcommand\footnoterule{\kern-3\p@ \hrule width 0.4\columnwidth \kern 2.6\p@}
\@addtoreset{footnote}{page}

\long\def\@makefntext#1{\@setpar{\@@par\@tempdima \hsize
\advance\@tempdima-\@footindent
\parshape \@ne \@footindent \@tempdima}\par
\noindent \hbox to \z@{\hss\@thefnmark\hspace{0.5em}}#1}

\renewcommand\thefootnote{\pinumber{\arabic{footnote}}}
\def\@makefnmark{\hbox{\textsuperscript{\@thefnmark}}}

\newcommand\pinumber[1]{\ifcase#1 \or \ding{172}\or \ding{173}\or
  \ding{174}\or \ding{175}\or \ding{176}\or \ding{177}%
  \or \ding{178}\or \ding{179}\or \ding{180}\or \ding{181}\else *\fi\relax}
\makeatother
%以上从\makeatletter到\makeatother为重定义脚注编号,使之带圆圈

\BiChapter{绪论}{Introduction}
\label{cha1:Introduction}

\BiSection{研究背景}{Background} %此处的第二个{}内是用于填入英语标题，南理工没要求，所以空着，但不要删除它，下同。
\label{sec11:Background}

XXXXXXXXXXXXXXXXXXXXXXXXXXXXXXXXXXXXXXXX

XXXXXXXXXXXXXXXXxxxx

XXXXXXxx \cite{psychology_recognition_87}



% !Mode:: "TeX:UTF-8"
% !TeX root = ../main.tex
% 这个是为了WinEdt设置的，它的默认不是UTF8.

% 若xelatex编译非UTF8文件，需在每个文件中指定字符编码;
% main.tex中手动制定了\atemp和\usewhat参数；
\ifx\atempxetex\usewhat
%\XeTeXinputencoding "gbk"
\fi
\defaultfont

\BiChapter{XXXXXX}{XXXXXx}
\label{cha2}

\BiSection{XXXX}{XXXXXXX}
\label{sec20}



\clearpage %----这个命令是什么意思？




% !Mode:: "TeX:UTF-8"
% !TeX root = ../main.tex
% 若xelatex编译非UTF8文件，需在每个文件中指定字符编码;
% main.tex中手动制定了\atemp和\usewhat参数；
\ifx\atempxetex\usewhat
%\XeTeXinputencoding "gbk"
\fi
\defaultfont


\BiChapter{XXXXXXXXXXXX}{XXXXXXXXXXXXx}
\label{cha3}

\BiSection{引言}{Introduction}
\label{sec31}






% !Mode:: "TeX:UTF-8"
% !TeX root = ../main.tex
% 若xelatex编译非UTF8文件，需在每个文件中指定字符编码;
% main.tex中手动制定了\atemp和\usewhat参数；
\ifx\atempxetex\usewhat
%\XeTeXinputencoding "gbk"
\fi
\defaultfont



\BiChapter{XXXXXX}{XXXXXXXXXXXx}
\label{cha4}

\BiSection{引言}{Introduction}
\label{sec41}





% !Mode:: "TeX:UTF-8"
% !TeX root = ../main.tex
% 这个是为了WinEdt设置的，它的默认不是UTF8.

% 若xelatex编译非UTF8文件，需在每个文件中指定字符编码;
% main.tex中手动制定了\atemp和\usewhat参数；
\ifx\atempxetex\usewhat
%\XeTeXinputencoding "gbk"
\fi
\defaultfont

\BiChapter{XXXXx}{XXXXXXXXXXXXX }
\label{chp6}

\BiSection{引言}{Introduction}
\label{sec6}


% !Mode:: "TeX:UTF-8"
% !TeX root = ../main.tex
% 若xelatex编译非UTF8文件，需在每个文件中指定字符编码;
% main.tex中手动制定了\atemp和\usewhat参数；
\ifx\atempxetex\usewhat
%\XeTeXinputencoding "gbk"
\fi
\defaultfont


\titleformat{\chapter}[hang]{\xiaosan\bf\filcenter\hei\sf\boldmath}{\xiaoer\chaptertitlename}{18pt}{\xiaosan}
\titlespacing{\chapter}{0pt}{8pt}{16pt}

\BiAppendixChapter{结~~~~束~~~~语}{Conclusion}


XXXXXXXXXXXXXXXXXXXXXXXXXXXXx

XXXXXXx

   % 结束语
% !Mode:: "TeX:UTF-8"
% 若xelatex编译非UTF8文件，需在每个文件中指定字符编码;
% main.tex中手动制定了\atemp和\usewhat参数；
\ifx\atempxetex\usewhat
%\XeTeXinputencoding "gbk"
\fi
\defaultfont

\BiAppendixChapter{致~~~~谢}{Acknowledgement}

XXXXXXXXXXXXXXXXXXXXXXXXXXXXX




% 致谢
%%% ******** 参考文献***************%%%
\defaultfont
\ifx\atempxetex\usewhat
\bibliographystyle{chinesebst2005}
\else
\bibliographystyle{chinesebst}
\fi
\addcontentsline{toc}{chapter}{\hei \ReferenceCName}      % 参考文献加入到中文目录
\addcontentsline{toe}{chapter}{\bfseries \ReferenceEName} % 参考文献加入到英文目录
\addtolength{\bibsep}{-0.8 em}

%\nocite{*}  %如果不注释掉，会将bib文件中所有的引文都列出来，不管是否在正文中引用。
\bibliography{reference/reference}

%\addtocontents{fen}{\protect\vskip1\baselineskip}
%\addtocontents{ten}{\protect\vskip1\baselineskip}
%英文图形和表格索引里加入空白行，通常放在 % !Mode:: "TeX:UTF-8"
% 这个是为了WinEdt设置的，它的默认不是UTF8.
% !TeX root = ../main.tex
% 若xelatex编译非UTF8文件，需在每个文件中指定字符编码;
% main.tex中手动制定了\atemp和\usewhat参数；
\ifx\atempxetex\usewhat
%\XeTeXinputencoding "gbk"
\fi



%%%%%%%%%%%%%%%%%%%%%%%%%%%%%%%%%%%%%%%%%%%%%%%%%%%%%%%%
 \defaultfont
 \appendix

 % %%%%%%%%%%%%%%%%%%%%%%%%%%%%%%%%%%%%%%%%%%%%%%%%%%%%%%%%%
% 附录A之前。
%区分开正文和附录的图形和表格，一般没有这个必要。

%%% ******** 附件***************%%%
%%------------注意(womegaga): 下面三行是撰写每一章时，为避免不必要的编译，所以注释掉，最终论文需要去掉下面三行注释，并进行相应的核实和修改%%

% !Mode:: "TeX:UTF-8"
% 这个是为了WinEdt设置的，它的默认不是UTF8.
% !TeX root = ../main.tex
% 若xelatex编译非UTF8文件，需在每个文件中指定字符编码;
% main.tex中手动制定了\atemp和\usewhat参数；
\ifx\atempxetex\usewhat
%\XeTeXinputencoding "gbk"
\fi



%%%%%%%%%%%%%%%%%%%%%%%%%%%%%%%%%%%%%%%%%%%%%%%%%%%%%%%%
 \defaultfont
 \appendix

 % %%%%%%%%%%%%%%%%%%%%%%%%%%%%%%%%%%%%%%%%%%%%%%%%%%%%%%%%%
            % 附录A
% !Mode:: "TeX:UTF-8"
% 若xelatex编译非UTF8文件，需在每个文件中指定字符编码;
% main.tex中手动制定了\atemp和\usewhat参数；
\ifx\atempxetex\usewhat
%\XeTeXinputencoding "gbk"
\fi
\defaultfont
\BiAppendixChapter{附~~~~录}{Fulu}
%\BiAppendixChapter{攻读\cxuewei 学位期间发表的学术论文} {Papers Published in the Period of PH. D. Education}

\song\sihao\sf\bf{攻读博士期间完成的已发表的论文：}
\vspace{10pt}
\begin{publist}

\xiaosi \item XXXXXX

\vspace{10pt}
\item XXXXXXX


\vspace{10pt}
\item XXXXXx





\end{publist}



\vspace{100pt}
\song\sihao\sf\bf{作者在攻读博士学位期间参与的主要科研项目：}
\vspace{10pt}
\begin{publist}
\xiaosi \item XXXXXXX。
 \vspace{10pt}
 \item XXXXXX”。
\end{publist}
    % 所发文章
% !Mode:: "TeX:UTF-8"
% 若xelatex编译非UTF8文件，需在每个文件中指定字符编码;
% main.tex中手动制定了\atemp和\usewhat参数；
%\ifx\atempxetex\usewhat
%%\XeTeXinputencoding "gbk"
%\fi
%\defaultfont
%
%\BiAppendixChapter{个人简历}{Resume}
%
%{\hei 学习经历}
%\begin{publist}
%\item 2007~年~3~月--至今~~南京理工大学计算机科学与技术学院~~~攻读工学博士学位
%\item 2004~年~9~月~~2006~年~12~月~~~南京林业大学信息技术学院~~~获工学硕士学位
%\item 2000~年~9~月--2004~年~6~月~~~南京邮电学院应用数理系~~~获理学学士学位
%\end{publist}
%
%{\hei 科研工作}
%\begin{publist}
%\item  2007~年~x~月--xxxx~年~x~月 ~~~ xxxx项目~~~~（编号xxx-xxx-xxx）　
%\item  xxxx~年~x~月--xxxx~年~x~月 ~~~ xxxx项目~~~~（编号xxx-xxx-xxx）
%\item  xxxx~年~x~月--xxxx~年~x~月 ~~~ xxxx项目~~~~（编号xxx-xxx-xxx）
%\item  xxxx~年~x~月--xxxx~年~x~月 ~~~ xxxx项目~~~~（编号xxx-xxx-xxx）
%\end{publist}
%
%{\hei 学术论文}
%\begin{publist}
%\item 在~xxxxxxx~等刊物发表论文多篇
%\item 在~xxxxxxxxxxxxxxxx~等多个国际会议上发表论文多篇
%\end{publist}
%
          % 个人简历

\clearpage
\ifx\atempxetex\usewhat\else
\end{CJK*}
\fi

\end{document}
